\section{Implementácia navrhnutého riešenia}

V tejto kapitole je opísaná implementácia základného riešenia pre odpovedanie
na otázky nad tabuľkovými dátami z datasetu DataBench. Implementované riešenie
využíva jazykový model Gemini na generovanie Python kódu, ktorý následne
vykonáva požadovanú operáciu nad tabuľkovými dátami. Výsledok vykonania je
normalizovaný a následne porovnaný s referenčnou odpoveďou pomocou nástroja
\texttt{databench\_eval}.

Implementácia je rozdelená do viacerých samostatných modulov. Každý modul
zodpovedá za konkrétnu časť spracovania, čo umožňuje jednoduchšie testovanie
a prípadné úpravy jednotlivých častí systému.

Projekt je rozdelený do niekoľkých Python súborov. Hlavný priebeh programu
zabezpečuje súbor \texttt{main.py}, ktorý spája jednotlivé komponenty do
jedného procesu. Načítanie tabuliek zabezpečuje modul
\texttt{dataset\_loader.py}, tvorbu promptu modul
\texttt{prompt\_builder.py} a komunikáciu s jazykovým modelom modul
\texttt{llm\_client.py}.

Vykonávanie vygenerovaného kódu je oddelené do modulov
\texttt{code\_executor.py} a \texttt{sandbox\_runner.py}. Modul
\texttt{evaluator.py} zabezpečuje porovnanie výslednej odpovede s referenčnou
odpoveďou a modul \texttt{metrics.py} vypočítava základné metriky z uložených
výsledkov.

\begin{verbatim}
DataBench
    |
    v
main.py
    |
    +--> dataset_loader.py
    |       |
    |       +--> all.parquet / sample.parquet
    |
    +--> prompt_builder.py
    |       |
    |       +--> vytvorenie promptu
    |
    +--> llm_client.py
    |       |
    |       +--> Gemini
    |       |
    |       +--> Python kód
    |
    +--> code_executor.py
    |       |
    |       +--> kontrola syntaxe
    |       +--> Docker
    |              |
    |              +--> sandbox_runner.py
    |
    +--> evaluator.py
    |       |
    |       +--> porovnanie odpovede
    |
    +--> results.jsonl
            |
            +--> metrics.py
\end{verbatim}

Takéto rozdelenie umožňuje oddeliť načítanie dát, generovanie kódu, jeho
vykonanie a vyhodnotenie výsledkov.

\subsection{Načítanie otázok a tabuľkových dát}

Hlavný program načítava otázky z datasetu DataBench pomocou knižnice
\texttt{datasets}. V aktuálnej implementácii sa používa konfigurácia
\texttt{qa} a split \texttt{train}. Každá otázka obsahuje okrem samotného
textu otázky aj identifikátor tabuľky, typ očakávanej odpovede a referenčnú
odpoveď.

Pre každú otázku sa podľa identifikátora datasetu vyhľadá príslušná lokálna
tabuľka vo formáte Parquet. Implementácia podporuje dva režimy spracovania.
V režime \texttt{full} sa používa súbor \texttt{all.parquet}, ktorý obsahuje
celú tabuľku. V režime \texttt{lite} sa používa súbor
\texttt{sample.parquet}, ktorý predstavuje menšiu vzorku dát.

Na načítanie tabuľky sa používa knižnica \texttt{pandas}. Po načítaní tabuľky
sa pomocou funkcie \texttt{get\_table\_schema()} vytvorí schéma obsahujúca
názvy jednotlivých stĺpcov a ich dátové typy. Táto informácia sa následne
používa pri tvorbe promptu pre jazykový model.

\subsection{Tvorba promptu a Generovanie Python kódu}

Po načítaní otázky a tabuľky sa vytvorí prompt pomocou funkcie
\texttt{build\_prompt()}. Prompt obsahuje samotnú otázku, zoznam dostupných
stĺpcov, ich dátové typy a očakávaný typ odpovede.

Dôležitou súčasťou promptu sú pravidlá pre generovanie výsledného Python
kódu. Jazykový model dostáva informáciu, že objekt \texttt{df} obsahujúci
tabuľku je už pripravený. Vygenerovaný program musí uložiť výsledok do
premennej \texttt{result}.

Prompt zároveň obsahuje pravidlá pre jednotlivé typy odpovedí. Podporované
sú napríklad booleovské hodnoty, čísla, reťazce, kategórie a zoznamy
kategórií alebo čísel. Týmto spôsobom sa znižuje riziko, že jazykový model
vráti textové vysvetlenie namiesto požadovanej hodnoty.

V prompte sú tiež uvedené pravidlá týkajúce sa spracovania duplicitných
hodnôt, zoradenia výsledkov, číselných hodnôt a konverzie dátových typov.
Model je zároveň upozornený, že nesmie používať stĺpce alebo hodnoty, ktoré
nie sú súčasťou poskytnutej tabuľky.


Na generovanie Python kódu sa používa model Gemini prostredníctvom knižnice
\texttt{google-genai}. API kľúč je načítaný z environmentálnej premennej
\texttt{GEMINI\_API\_KEY}, ktorá je uložená v súbore \texttt{.env}.

Model dostane pripravený prompt a jeho úlohou je vytvoriť vykonateľný Python
kód. Výstup modelu je následne upravený pomocou funkcie
\texttt{clean\_generated\_code()}, ktorá odstraňuje prípadné Markdown
ohraničenie kódu, napríklad bloky označené pomocou \texttt{```python}.

Implementácia obsahuje aj opakovanie požiadavky pri dočasných chybách
komunikácie s API. Pri chybách súvisiacich s prekročením limitu požiadaviek
sa používa postupné predlžovanie času čakania medzi jednotlivými pokusmi.
Maximálny počet pokusov komunikácie s modelom je nastavený na päť.

Pred samotným vykonaním kódu sa kontroluje jeho syntaktická správnosť.
Na tento účel sa používa modul \texttt{ast} a funkcia
\texttt{ast.parse()}.

Ak kód obsahuje syntaktickú chybu, nie je odoslaný na vykonanie. Systém
zaznamená chybu a môže požiadať jazykový model o novú verziu kódu. V rámci
spracovania jednej otázky sú povolené maximálne dve generovania kódu --
prvotný pokus a jedna opakovaná generácia.

Pri opakovanom pokuse je model informovaný o tom, že predchádzajúci
vygenerovaný kód nebol úspešný. Do promptu sa vloží predchádzajúci kód
a informácia o chybe, aby mohol model vytvoriť opravenú verziu.

\subsection{Izolované vykonávanie Python kódu}

Vygenerovaný Python kód sa nevykonáva priamo v hlavnom procese programu.
Na jeho vykonanie sa používa Docker kontajner. Tento prístup umožňuje oddeliť
vykonávanie kódu generovaného jazykovým modelom od hlavnej aplikácie.

Kontajner sa spúšťa s viacerými obmedzeniami. Sieťová komunikácia je
zakázaná pomocou parametra \texttt{--network none}. Kontajner má zároveň
obmedzenú operačnú pamäť na 512 MB a využitie procesora na jeden CPU.
Počet procesov je obmedzený pomocou parametra \texttt{--pids-limit}.

Súborový systém kontajnera je nastavený ako iba na čítanie. Pre dočasné
operácie je povolený samostatný adresár \texttt{/tmp}. Zároveň sú odstránené
Linuxové capabilities pomocou \texttt{--cap-drop ALL} a je použitá možnosť
\texttt{no-new-privileges}.

Do kontajnera sa pripájajú iba potrebné súbory. Ide o príslušnú tabuľku
\texttt{table.parquet}, súbor so vygenerovaným kódom
\texttt{generated\_code.py} a súbor \texttt{sandbox\_runner.py}. Všetky
tieto súbory sú pripojené v režime iba na čítanie.

Samotné vykonanie je časovo obmedzené na 20 sekúnd. Ak program prekročí
tento limit alebo Docker vráti chybu, vykonanie sa považuje za neúspešné.

\subsection{Spracovanie výsledku v sandboxe}

Po spustení kontajnera načíta súbor \texttt{sandbox\_runner.py} príslušnú
tabuľku do objektu \texttt{df}. Vygenerovaný Python kód sa následne vykoná
v priestore, v ktorom je objekt \texttt{df} dostupný.

Po vykonaní kódu sa kontroluje, či vygenerovaný program vytvoril premennú
\texttt{result}. Ak táto premenná neexistuje, vykonanie sa považuje za
neúspešné.

Výsledok sa následne normalizuje pomocou funkcie
\texttt{normalize\_result()}. Táto funkcia konvertuje hodnoty z knižníc
Pandas a NumPy na štandardné Python dátové typy, ktoré je možné serializovať
do formátu JSON. Podporované sú napríklad skalárne hodnoty, zoznamy,
slovníky a hodnoty získané z objektov \texttt{Series} alebo
\texttt{ndarray}.

Takto pripravený výsledok je následne odoslaný späť do hlavného programu
vo forme JSON. Hlavný program ho dekóduje a používa ako predikovanú odpoveď
na danú otázku.

\subsection{Opakovanie po chybe vykonávania}

Ak sa vygenerovaný kód nepodarí vykonať v Docker kontajneri, systém sa
pokúsi vytvoriť nový kód. Maximálne sú teda použité dva pokusy v rámci
spracovania jednej otázky.

Pri druhom pokuse sa do promptu pridá informácia o predchádzajúcom
neúspešnom vykonaní a predchádzajúci vygenerovaný kód. Jazykový model je
požiadaný o vytvorenie opravenej verzie.

Aktuálna implementácia používa pri chybe Dockeru všeobecnú informáciu
o neúspešnom vykonaní. Konkrétny traceback z Docker kontajnera sa síce
vypíše do konzoly, ale ešte nie je štruktúrovane odovzdaný späť jazykovému
modelu.

\subsection{Vyhodnotenie odpovede}

Po úspešnom vykonaní Python kódu sa získaná hodnota porovná s referenčnou
odpoveďou. Na tento účel sa používa trieda \texttt{Evaluator} z balíka
\texttt{databench\_eval}.

Funkcia \texttt{evaluate\_answer()} využíva porovnávaciu funkciu
\texttt{default\_compare()}, ktorá zohľadňuje aj typ odpovede. Výsledkom
porovnania je logická hodnota určujúca, či sa predikovaná odpoveď považuje
za správnu.

V režime \texttt{full} sa ako referenčná odpoveď používa hodnota
\texttt{answer}. V režime \texttt{lite} sa používa hodnota
\texttt{sample\_answer}. V prípade, že referenčná odpoveď chýba, otázka
sa označí stavom \texttt{unscored} a nepovažuje sa za úspešne vyhodnotenú.

\subsection{Ukladanie výsledkov}

Výsledky jednotlivých otázok sa ukladajú do súboru
\texttt{results.jsonl}. Každý riadok obsahuje jeden JSON objekt.

Medzi ukladané informácie patria napríklad index otázky, identifikátor
tabuľky, režim spracovania, samotná otázka, typ odpovede, očakávaná
odpoveď, predikovaná odpoveď, informácia o správnosti, stav spracovania,
chybová správa, vygenerovaný kód a počet pokusov.

Takéto uloženie umožňuje neskoršie analyzovať jednotlivé výsledky bez
potreby opätovného spúšťania jazykového modelu.

\subsection{Režimy spustenia programu}

Program \texttt{main.py} podporuje dva režimy spracovania otázok:
\texttt{single} a \texttt{batch}.

Režim \texttt{single} umožňuje vybrať jednu otázku podľa jej indexu.
Tento režim je vhodný najmä na manuálne overovanie správania systému,
promptu a generovaného kódu.

Režim \texttt{batch} umožňuje spracovať súvislý rozsah otázok. Používateľ
zadá počiatočný index a počet otázok. Program následne postupne spracuje
jednotlivé otázky a priebežne ukladá ich výsledky do súboru
\texttt{results.jsonl}.

Oba režimy podporujú režim dát \texttt{full} a \texttt{lite}. Týmto
spôsobom je možné experimentálne porovnať spracovanie úplných tabuliek
a menších vzoriek.

\subsection{Testovanie implementácie}

Súčasťou projektu sú aj testovacie skripty, ktoré overujú jednotlivé časti
implementácie. Testovanie zahŕňa napríklad kontrolu načítania lokálnych
tabuliek, správanie evaluátora, spracovanie chýb Docker kontajnera a
správanie opakovania po chybe.

Pre kontrolu syntaktickej správnosti Python súborov je možné použiť
modul \texttt{py\_compile}. Samostatný test
\texttt{test\_retry.py} je implementovaný ako štandardný
\texttt{unittest} test pre kontrolu opakovania po neúspešnom vykonaní.
Ostatné testovacie súbory sú v aktuálnej verzii projektu pripravené
predovšetkým ako manuálne validačné skripty.

Implementované riešenie realizuje základný baseline založený na generovaní
Python kódu jazykovým modelom. Celý proces pozostáva z načítania otázky
a tabuľky, vytvorenia schémy a promptu, generovania Python kódu, kontroly
jeho syntaxe, izolovaného vykonania v Docker kontajneri, normalizácie
výsledku a jeho porovnania s referenčnou odpoveďou.

Navrhnutá architektúra zároveň umožňuje zaznamenávať výsledky a chyby
jednotlivých otázok. To poskytuje základ pre experimentálne vyhodnotenie
baseline a pre ďalšie rozšírenia projektu v nasledujúcej etape.